% ============================================================
\chapter{Version Control --- Git and DVC}
\label{ch:version-control}
\index{version control}
% ============================================================

\section{Why Version Everything?}
\label{sec:why-version}

In software engineering, versioning code with Git is standard practice.
In ML, we must also version \textbf{data} and \textbf{models}---two
artefacts that evolve independently of the code and that jointly determine
the behaviour of the system.

\begin{definition}[Reproducibility]
An ML experiment is \textbf{reproducible}\index{reproducibility} if, given
the same code, the same data, and the same configuration, it produces the
same model and the same metrics (up to hardware non-determinism).
\end{definition}

\begin{remark}
Git alone is insufficient for ML projects: datasets and model weights are
often too large for Git (files $>100$\,MB are impractical), and binary
files do not benefit from line-level diffing.
\end{remark}

\section{Git Essentials for ML}
\label{sec:git-essentials}
\index{Git}

\subsection{Core workflow}

\begin{codeblock}[title=\textbf{Essential Git commands for ML projects}]
\begin{minted}{bash}
# Initialise a repository
git init

# Stage and commit
git add src/ configs/
git commit -m "Add training pipeline"

# Create a feature branch
git checkout -b experiment/new-feature-engineering

# Merge back
git checkout main
git merge experiment/new-feature-engineering

# Tag a release (e.g., after a successful model)
git tag -a v1.0.0 -m "Baseline model, accuracy 0.87"
\end{minted}
\end{codeblock}

\subsection{.gitignore for ML projects}
\index{.gitignore}

\begin{codeblock}[title=\textbf{.gitignore tailored for ML}]
\begin{minted}{text}
# Data and models (managed by DVC)
data/
models/
*.pkl
*.joblib
*.h5
*.pt
*.onnx

# Environments
.venv/
__pycache__/
*.pyc

# Notebooks
.ipynb_checkpoints/

# Experiment artefacts
mlruns/
wandb/
outputs/

# Secrets
.env
*.secret
\end{minted}
\end{codeblock}

\begin{bestpractice}[title=\textbf{Git branching strategy for ML}]
\begin{itemize}
  \item \textbf{main}: production-ready code and pipeline definitions.
  \item \textbf{develop}: integration branch for ongoing work.
  \item \textbf{experiment/*}: one branch per experiment (short-lived).
  \item \textbf{data/*}: branches for data processing changes.
  \item Use \textbf{tags} to mark successful model versions
    (e.g., \cli{v1.2.0-acc0.91}).
\end{itemize}
\end{bestpractice}

\section{DVC --- Data Version Control}
\label{sec:dvc}
\index{DVC}

\textbf{DVC}\index{DVC} (Data Version Control) extends Git to handle large
files (datasets, model weights) without storing them in the Git repository.
Instead, DVC stores metadata (\file{.dvc} files) in Git and the actual data
in a configurable remote storage (S3, GCS, Azure Blob, local, SSH).

\begin{definition}[DVC]
DVC is an open-source version control system for data and ML models.
It works \emph{alongside} Git: Git tracks code and \file{.dvc} metadata
files, while DVC tracks the actual large files via content-addressable
storage.
\end{definition}

\subsection{Getting started with DVC}

\begin{codeblock}[title=\textbf{DVC initialisation and basic usage}]
\begin{minted}{bash}
# Install DVC
pip install dvc dvc-s3  # add dvc-gdrive, dvc-azure, etc.

# Initialise DVC in an existing Git repo
cd my-ml-project
dvc init

# Track a dataset
dvc add data/train.csv
# Creates: data/train.csv.dvc  (metadata, tracked by Git)
#          data/.gitignore      (so Git ignores train.csv)

git add data/train.csv.dvc data/.gitignore
git commit -m "Track training data with DVC"

# Configure remote storage
dvc remote add -d myremote s3://my-bucket/dvc-store

# Push data to remote
dvc push

# Pull data on another machine
git clone <repo-url>
dvc pull
\end{minted}
\end{codeblock}

\subsection{.dvc file anatomy}

\begin{codeblock}[title=\textbf{data/train.csv.dvc}]
\begin{minted}{yaml}
outs:
  - md5: a1b2c3d4e5f6a1b2c3d4e5f6a1b2c3d4
    size: 157286400
    hash: md5
    path: train.csv
\end{minted}
\end{codeblock}

The \file{.dvc} file records the MD5 hash and size of the tracked file.
When you run \cli{dvc push}, DVC uploads the file to the remote, indexed
by its hash. When you run \cli{dvc pull}, DVC downloads the file matching
the hash in the \file{.dvc} file.

\subsection{Switching between data versions}

\begin{codeblock}[title=\textbf{Time-travelling through data versions}]
\begin{minted}{bash}
# Check out a previous data version
git checkout v1.0 -- data/train.csv.dvc
dvc checkout

# Compare data versions
dvc diff HEAD~1

# List tracked files
dvc list . --dvc-only
\end{minted}
\end{codeblock}

\subsection{.dvcignore}
\index{.dvcignore}

Like \file{.gitignore}, the \file{.dvcignore} file tells DVC which files
to exclude when scanning directories.

\begin{codeblock}[title=\textbf{.dvcignore}]
\begin{minted}{text}
# Temporary files
*.tmp
*.log

# OS files
.DS_Store
Thumbs.db

# Intermediate artefacts
data/raw/scratch/
\end{minted}
\end{codeblock}

\section{DVC Pipelines}
\label{sec:dvc-pipelines}
\index{DVC!pipelines}

DVC can define reproducible ML pipelines in a \file{dvc.yaml} file.
Each stage declares its dependencies and outputs; DVC tracks which stages
need to be re-run when inputs change.

\begin{codeblock}[title=\textbf{dvc.yaml --- Complete ML pipeline}]
\begin{minted}{yaml}
stages:
  prepare:
    cmd: python src/prepare.py --config configs/config.yaml
    deps:
      - src/prepare.py
      - data/raw/
      - configs/config.yaml
    outs:
      - data/processed/train.csv
      - data/processed/test.csv

  train:
    cmd: python src/train.py --config configs/config.yaml
    deps:
      - src/train.py
      - data/processed/train.csv
      - configs/config.yaml
    outs:
      - models/model.pkl
    metrics:
      - metrics/train_metrics.json:
          cache: false

  evaluate:
    cmd: python src/evaluate.py --config configs/config.yaml
    deps:
      - src/evaluate.py
      - models/model.pkl
      - data/processed/test.csv
    metrics:
      - metrics/eval_metrics.json:
          cache: false
    plots:
      - metrics/confusion_matrix.csv:
          x: predicted
          y: actual
\end{minted}
\end{codeblock}

\begin{codeblock}[title=\textbf{Running and inspecting DVC pipelines}]
\begin{minted}{bash}
# Run the full pipeline (only re-runs changed stages)
dvc repro

# Visualise the DAG
dvc dag

# Compare metrics across Git commits
dvc metrics diff

# Show metrics
dvc metrics show

# Show plots
dvc plots show
\end{minted}
\end{codeblock}

\begin{remark}
DVC uses file hashes to determine whether a stage needs re-execution.
If neither the dependencies nor the code have changed, the stage is
skipped---making iteration fast.
\end{remark}

\section{Git + DVC Workflow}
\label{sec:git-dvc-workflow}

\begin{center}
\begin{tikzpicture}[
  box/.style={rectangle, draw, rounded corners, minimum width=3.2cm,
    minimum height=1.4cm, align=center, font=\small},
  storage/.style={cylinder, draw, shape border rotate=90,
    minimum width=2cm, minimum height=1.2cm, align=center,
    font=\scriptsize, aspect=0.3},
  arr/.style={-{Stealth[length=2.5mm]}, thick},
  lbl/.style={font=\scriptsize, midway, fill=white, inner sep=1pt}
]
  % Developer
  \node[box, fill=blue!10] (dev) {\textbf{Developer}\\Working Directory};

  % Git
  \node[storage, fill=green!10, above right=2cm and 2.5cm of dev] (git)
    {\textbf{Git Remote}\\Code, configs,\\\file{.dvc} files};

  % DVC
  \node[storage, fill=orange!10, below right=2cm and 2.5cm of dev] (dvc)
    {\textbf{DVC Remote}\\Data, models\\(S3, GCS, \ldots)};

  % Arrows
  \draw[arr, steelblue] (dev.north east) -- (git.west)
    node[lbl, above, sloped] {\cli{git push}};
  \draw[arr, steelblue] (git.west) -- (dev.north east)
    node[lbl, below, sloped] {\cli{git pull}};

  \draw[arr, deepred] (dev.south east) -- (dvc.west)
    node[lbl, above, sloped] {\cli{dvc push}};
  \draw[arr, deepred] (dvc.west) -- (dev.south east)
    node[lbl, below, sloped] {\cli{dvc pull}};

  % Labels
  \node[font=\scriptsize, text=steelblue, above=0.3cm of git]
    {Code + metadata};
  \node[font=\scriptsize, text=deepred, below=0.3cm of dvc]
    {Large files};
\end{tikzpicture}
\end{center}

\section{Tool Comparison}
\label{sec:version-comparison}

\begin{center}
\begin{tabular}{@{}lcccc@{}}
\toprule
\textbf{Criterion} & \textbf{DVC} & \textbf{Git LFS} & \textbf{Delta Lake} & \textbf{LakeFS} \\
\midrule
Open source        & Yes  & Yes  & Yes  & Yes  \\
Git integration    & Excellent & Native & No & Git-like \\
Pipeline support   & Yes  & No   & No   & No   \\
Storage backends   & Many & Git server & Object store & S3-compatible \\
Data format        & Any  & Any  & Parquet & Any  \\
Branching for data & Via Git & Via Git & No & Yes \\
Scalability        & Good & Limited & Excellent & Excellent \\
Learning curve     & Low  & Very low & Medium & Medium \\
\bottomrule
\end{tabular}
\end{center}

\begin{warning}[title=\textbf{Git LFS limitations}]
Git LFS stores large files on the Git server itself, which can become
expensive and slow for multi-GB datasets. DVC decouples data storage from
the Git server, allowing you to use cheap cloud storage (S3, GCS).
Git LFS also lacks pipeline support and metric tracking.
\end{warning}

\section{Best Practices}

\begin{bestpractice}[title=\textbf{Version control golden rules}]
\begin{enumerate}
  \item \textbf{Commit often, commit small}: each commit should represent
    one logical change.
  \item \textbf{Never commit data to Git}: use DVC or a similar tool.
  \item \textbf{Never commit secrets}: use \file{.env} files excluded via
    \file{.gitignore}.
  \item \textbf{Tag successful experiments}: \cli{git tag v1.0-acc0.92}.
  \item \textbf{Use DVC pipelines}: make your workflow reproducible with
    \file{dvc.yaml}.
  \item \textbf{Automate with hooks}: pre-commit hooks for linting,
    formatting, secret scanning.
\end{enumerate}
\end{bestpractice}

\begin{antipattern}[title=\textbf{Common anti-patterns}]
\begin{itemize}
  \item Storing datasets in Git (slow clone, bloated history)
  \item Using \cli{model\_v2\_final\_FINAL.pkl} naming instead of proper
    versioning
  \item Not versioning data at all (``I'll remember which file I used'')
  \item Giant commits with thousands of changed files
  \item Forgetting to \cli{dvc push} after \cli{git push}
\end{itemize}
\end{antipattern}

\section{Mini-project: Versioned ML Pipeline}
\label{sec:mini-project-ch3}

\begin{miniproject}[title=\textbf{Mini-project 3: Git + DVC pipeline}]
Building on the project from Chapters~\ref{ch:introduction}
and~\ref{ch:environments}:
\begin{enumerate}
  \item Initialise a Git repository and DVC.
  \item Track a dataset ($\ge 50$\,MB) with DVC and configure a remote
    (local directory or S3).
  \item Write a \file{dvc.yaml} with three stages: \texttt{prepare},
    \texttt{train}, \texttt{evaluate}.
  \item Run \cli{dvc repro} and verify that only changed stages re-run.
  \item Tag the first successful run as \cli{v1.0}.
  \item Modify a hyperparameter, re-run, and use \cli{dvc metrics diff}
    to compare with the previous version.
  \item Push code to Git and data to DVC remote; verify a colleague can
    reproduce the results with \cli{git clone} + \cli{dvc pull} +
    \cli{dvc repro}.
\end{enumerate}
\end{miniproject}

\section{Exercises}
\label{sec:exercises-ch3}

\begin{exercise}[$\bigstar$ --- DVC basics]
Initialise a Git+DVC project. Track a CSV file ($\ge 10$\,MB) with DVC.
Configure a local remote (\cli{dvc remote add -d local /tmp/dvc-store}).
Push and pull the file. Verify the \file{.dvc} file contents and explain
each field.
\end{exercise}

\begin{exercise}[$\bigstar\bigstar$ --- Pipeline design]
Design and implement a DVC pipeline (\file{dvc.yaml}) for a text
classification task with the following stages:
\begin{enumerate}
  \item \texttt{download}: download the dataset from a URL
  \item \texttt{preprocess}: tokenisation, vocabulary building
  \item \texttt{train}: train a model, log metrics to JSON
  \item \texttt{evaluate}: compute precision, recall, F1 on the test set
\end{enumerate}
Run the pipeline, change a hyperparameter, re-run, and compare metrics
with \cli{dvc metrics diff}.
\end{exercise}

\begin{exercise}[$\bigstar\bigstar\bigstar$ --- Collaborative workflow]
Simulate a two-person collaboration using Git branches and DVC:
\begin{enumerate}
  \item Person A creates an experiment branch, modifies the feature
    engineering, trains, and pushes code + data
  \item Person B creates a different experiment branch with a different
    model architecture
  \item Both merge into \texttt{develop}, resolving any conflicts in
    \file{dvc.yaml} and \file{dvc.lock}
  \item Compare results using \cli{dvc metrics diff} across branches
  \item Write a short report documenting the collaboration workflow and
    any issues encountered
\end{enumerate}
\end{exercise}

\section{Cheatsheet}

\begin{cheatsheet}[title=\textbf{Chapter 3 Summary}]
\begin{tabular}{@{}p{4cm}p{8.5cm}@{}}
\toprule
\textbf{Tool / Concept} & \textbf{Key Commands} \\
\midrule
\textbf{Git basics} & \cli{git init}, \cli{git add}, \cli{git commit},
  \cli{git tag} \\
\textbf{DVC init} & \cli{dvc init}, \cli{dvc add <file>},
  \cli{dvc remote add} \\
\textbf{DVC data} & \cli{dvc push}, \cli{dvc pull}, \cli{dvc checkout},
  \cli{dvc diff} \\
\textbf{DVC pipelines} & \file{dvc.yaml}, \cli{dvc repro}, \cli{dvc dag},
  \cli{dvc metrics show} \\
\textbf{Best practice} & Git for code + \file{.dvc} files; DVC remote for
  data + models \\
\bottomrule
\end{tabular}
\end{cheatsheet}
